\subsection{习题}\label{sec:2.10}
\begin{exercise}\label{exer:2.1}
本题说明，\autoref{prop:2.1}中的 $\sqrt d$ 因子无法改进.
\begin{enumerate}
\item[(a)]\label{item:2.1a} 证明存在正的普适常数 $c$，使对任意 $n\ge1$，$\mathbb R^d$ 中半径为 $c\sqrt d\,n^{-1/d}$ 的球，其 Lebesgue 测度至多为 $(2n)^{-1}$.
\par 提示：回顾 $\mathbb R^d$ 的单位球的 Lebesgue 测度为
$$\pi^{d/2}/\Gamma(d/2+1).$$
\item[(b)] 令 $\mu$ 是 $[0,1]^d$ 上的均匀测度，令 $\tilde\mu_n$ 是任意支撑在 $n$ 个点 $x_1,\ldots,x_n$ 上的测度. 用 $B(x_i,\epsilon)$ 表示以 $x_i$ 为中心、半径 $\epsilon$ 的球. 证明当 $\epsilon=c\sqrt d\,n^{-1/d}$ 时，$\mu(\bigcup_{i=1}^nB(x_i,\epsilon))\le1/2$，其中 $c$ 为\hyperref[item:2.1a]{第 (a) 问}的常数.
\item[(c)] 推出，若 $\gamma\in\Gamma(\mu,\tilde\mu_n)$，则 $\int\lVert x-y\rVert\,d\gamma(x,y)\ge\frac12\cdot c\sqrt d\,n^{-1/d}$.
\end{enumerate}
\end{exercise}

\begin{exercise}\label{exer:2.2}
修改\autoref{prop:2.5}的证明，建立 \eqref{eq:2.21}.
\end{exercise}

\begin{exercise}\label{exer:2.3}
回顾有界差分不等式（bounded differences inequality；例如\cite[定理6.2]{ref51}）. 利用它为 $W_1(\mu_n,\mu)$ 绕其期望的偏离证明一个集中不等式；这里 $X_1,\ldots,X_n$ 独立同分布于支撑包含在半径为 $R$ 的球中的分布 $\mu$，$\mu_n$ 为经验测度 $\mu_n=\frac1n\sum_{i=1}^n\delta_{X_i}$.
\end{exercise}

\begin{exercise}\label{exer:2.4}
证明，对任意 $p\ge1$，有 $\mathsf{SW}_p\le W_p$. 这个不等式是紧的吗？类似地，证明 $W_1^{(\sigma)}\le W_1$. 这个不等式是紧的吗？
\end{exercise}

\begin{exercise}\label{exer:2.5}
考虑\autoref{sec:2.8.4}定义的切片 Wasserstein 距离的计算问题.
\begin{enumerate}
\item[(a)] 证明，若 $\mu,\nu\in\mathcal P_p(\mathbb R)$（$p\ge1$），且两者分别均匀分布在 $n$ 个点上，则可以在 $O(n\log n)$ 时间内计算 $W_p(\mu,\nu)$；这里把算术和比较运算视为常数时间操作. 假设测度 $\mu,\nu$ 以未排序的点列表给出.
\item[(b)] 现在设 $\mu,\nu\in\mathcal P_p(\mathbb R^d)$ 分别均匀分布于 $n$ 个点，并以 $\mathbb R^d$ 中的点列表给出. 论证可在 $O(dn+n\log n)$ 时间内计算 $W_p(\Pi^\theta_\#\mu,\Pi^\theta_\#\nu)$.
\item[(c)] 从算法角度看这仍不够，因为精确计算 $\mathsf{SW}_1(\mu,\nu)$，需要对不可数多个 $\theta$ 计算 $W_1(\Pi^\theta_\#\mu,\Pi^\theta_\#\nu)$ 并积分. 转而论证：若 $\mu,\nu\in\mathcal P(B_1)$，从 $S^{d-1}$ 上的均匀测度 $\sigma$ 中独立同分布抽取 $m$ 个点 $\theta_1,\ldots,\theta_m$，则 Monte Carlo 平均
$$\widehat{\mathsf{SW}}_1(\mu,\nu):=\frac1m\sum_{i=1}^m W_1(\Pi^{\theta_i}_\#\mu,\Pi^{\theta_i}_\#\nu)$$
以 $O(m^{-1/2})$ 大小的加性误差逼近 $\mathsf{SW}_1(\mu,\nu)$.
\end{enumerate}
\end{exercise}

\begin{exercise}\label{exer:2.6}
\autoref{thm:1.7}说明，若 $\mu\in\mathcal P_1(\mathbb R)$，则
$$W_1(\mu_n,\mu)=\int_{-\infty}^{\infty}|F_{\mu_n}(t)-F_\mu(t)|\,dt,$$
其中 $F_{\mu_n}$ 和 $F_\mu$ 分别是 $\mu_n$ 和 $\mu$ 的累积分布函数. 利用这一事实，直接证明 $d=1$ 时的\autoref{prop:2.1}. 提示：$\mathbb E|F_{\mu_n}(t)-F_\mu(t)|^2=n^{-1}F_\mu(t)(1-F_\mu(t))$.
\end{exercise}

\begin{exercise}\label{exer:2.7}
\autoref{thm:2.15}证明的极小极大下界在 $d=1$ 时并非最优. 本题利用两个假设之间的检验（\cite[定理2.2]{ref302}）证明一个最优界. 使用这种方法，只需构造两个支撑包含于 $[0,1]$ 的测度 $\mu^{(0)},\mu^{(1)}$，使
$$\begin{aligned}
W_1(\mu^{(0)},\mu^{(1)})&\ge2r_n,\\
\mathrm{KL}(\mu^{(1)}\mathbin\Vert\mu^{(0)})&\le\frac1{2n}.
\end{aligned}$$
存在这样的测度意味着，对任何基于 $n$ 个独立同分布观测的估计量 $\tilde\mu_n$，存在 $j\in\{0,1\}$，使 $\mathbb E_{\mu^{(j)}}[W_1(\tilde\mu_n,\mu^{(j)})]\ge r_n/4$.
\begin{enumerate}
\item[(a)] 固定 $\varepsilon\in(0,1/2)$，考虑 $\mu^{(0)}=(\frac12+\varepsilon)\delta_0+(\frac12-\varepsilon)\delta_1$ 与 $\mu^{(1)}=(\frac12-\varepsilon)\delta_0+(\frac12+\varepsilon)\delta_1$. 证明 $W_1(\mu^{(0)},\mu^{(1)})=2\varepsilon$.
\item[(b)] 证明这对测度满足 $\mathrm{KL}(\mu^{(1)},\mu^{(0)})\le\frac{16\varepsilon^2}{1-4\varepsilon^2}$.
\item[(c)] 推出存在正的普适常数 $c$，使对任意估计量 $\tilde\mu_n$，存在 $j\in\{0,1\}$，满足
$$\mathbb E_{\mu^{(j)}}[W_1(\tilde\mu_n,\mu^{(j)})]\ge cn^{-1/2}.$$
\end{enumerate}
\end{exercise}

\begin{exercise}\label{exer:2.8}
\autoref{sec:2.8}讨论的正则化策略可以更广泛地使用. 例如，给定概率测度 $\mu,\nu$，定义如下量，并称为“平滑 $L^2$”：
$$d_\sigma^2(\mu,\nu):=\int\bigl(\mu\star\mathcal N(0,\sigma^2I)-\nu\star\mathcal N(0,\sigma^2I)\bigr)^2.$$
证明，它可以按参数速率估计：若 $\mu_n$ 表示由 $\mu$ 的 $n$ 个独立同分布样本形成的经验测度，则
$$\mathbb E d_\sigma^2(\mu_n,\mu)\le\frac1{(2\pi\sigma^2)^{d/2}n}.$$
\end{exercise}

\begin{exercise}\label{exer:2.9}
对于 $\mu,\nu\in\mathcal P_p(\mathbb R^d)$，它们之间的“最大切片 Wasserstein 距离”（max-sliced Wasserstein distance）\cite{ref109,ref182}\footnote{也称“低维 Wasserstein 距离”（low-dimensional Wasserstein distance）\cite{ref241}或“子空间稳健 Wasserstein 距离”（subspace robust Wasserstein distance）\cite{ref256}.} 为
$$\mathsf{MSW}_p(\mu,\nu):=\max_{\theta\in S}W_p(\Pi^\theta_\#\mu,\Pi^\theta_\#\nu).$$
本题目标是说明，$\mathsf{MSW}_1$ 可以按参数速率估计.
\begin{enumerate}
\item[(a)] 设 $\mu\in\mathcal P(B_1)$，$\mu_n$ 为 $X_1,\ldots,X_n\overset{\mathrm{i.i.d.}}\sim\mu$ 对应的经验测度. 证明
$$\mathsf{MSW}_1(\mu_n,\mu)=\sup_{f\in\mathcal F}\frac1n\sum_{i=1}^n\{f(X_i)-\mathbb Ef(X_i)\},$$
其中 $\mathcal F$ 是形如 $x\mapsto h(\theta^{\mathsf T}x)$ 的函数类，$\theta\in S^{d-1}$，$h$ 是 $[-1,1]$ 上满足 $h(0)=0$ 的 $1$-Lipschitz 函数.
\item[(b)] 证明
$$\log N(\varepsilon,\mathcal F)\lesssim1/\varepsilon+d\log(1+1/\varepsilon).$$
提示：参阅\autoref{sec:3}的\autoref{exer:3.2}.
\item[(c)] 利用\autoref{prop:2.6}，推出
$$\mathbb E\mathsf{MSW}_1(\mu_n,\mu)\lesssim\sqrt{d/n}.$$
事实上，更精细的论证表明，正确速率与维数无关\cite{ref44}.
\end{enumerate}
\end{exercise}
